\begin{table}[htpb]
\centering
\small
\caption{Resumen territorial y regional del recurso solar y potencial fotovoltaico en Córdoba (Atlas Solar Preliminar 2 km: Años Completos 2020, 2021 y 2023, Convenio ATLAS II).}
\label{tab:fase5b_metricas_provinciales}
\begin{tabular}{lcccccccc}
\toprule
\textbf{Región ($k=3$)} & \textbf{Celdas} & \textbf{GHI} & \textbf{DNI} & \textbf{$Y_{f,\mathrm{Fijo}}$} & \textbf{$Y_{f,\mathrm{Trk}}$} & \textbf{CF Fijo} & \textbf{CF Trk} & \textbf{Ganancia} \\
 & & \small{(kWh/m$^2$·a)} & \small{(kWh/m$^2$·a)} & \small{(kWh/kWp·a)} & \small{(kWh/kWp·a)} & \small{(\%)} & \small{(\%)} & \small{(\%)} \\
\midrule
Región 0 (Sierras Altas) & 5,586 & 1891.7 ± 69.3 & 2031.9 ± 119.4 & 1657.8 & 1993.3 & 18.9\% & 22.8\% & +20.22\% \\
Región 1 (Sierras Bajas / Trans.) & 17,677 & 1917.3 ± 106.3 & 2033.8 ± 228.5 & 1670.5 & 2014.9 & 19.1\% & 23.0\% & +20.59\% \\
Región 2 (Llanura Pampeana) & 25,382 & 1855.5 ± 37.4 & 1938.9 ± 58.1 & 1619.1 & 1943.9 & 18.5\% & 22.2\% & +20.05\% \\
\midrule
\textbf{Total Provincial} & \textbf{48,645} & \textbf{1882.1 ± 78.8} & \textbf{1984.0 ± 156.8} & \textbf{1642.2} & \textbf{1975.3} & \textbf{18.7\%} & \textbf{22.5\%} & \textbf{+20.27\%} \\
\bottomrule
\end{tabular}
\end{table}
